\documentclass[11pt,a4paper]{article}
\usepackage[T1]{fontenc}
\usepackage[utf8]{inputenc}
\usepackage[english]{babel}
\usepackage{lmodern}
\usepackage{amsmath,amssymb,amsthm}
\usepackage[top=2cm,bottom=2cm,left=2.5cm,right=2.5cm]{geometry}

\newtheorem{theorem}{Theorem}
\newtheorem{lemma}{Lemma}

\title{Improved Synchronous Leader Election\\on an Anonymous Ring}
\author{Mateusz Wojaczek}
\date{November 2026}

\begin{document}
\maketitle

\section{Introduction}

We consider a directed ring of $n$ anonymous processors: they run the same program, have no
identifiers, and processor $v_i$ can send messages only to $v_{i+1}$ (indices modulo $n$). The
network is synchronous, every processor knows $n$ and has its own random number generator. The
goal is to elect a leader, i.e.\ exactly one processor, so that at the end every processor knows
whether it is the leader. We present the improved algorithm of Itai and Rodeh~\cite{ir}
(Section~2.4).

\section{The algorithm}

The algorithm proceeds in phases of $n$ rounds, and all processors know a parameter $q$, $0<q<n$.
A \emph{pebble} is a message without content; it travels around the ring and is dropped by its
sender after exactly $n$ rounds. In every phase:
\begin{enumerate}
  \item each processor sends a pebble with probability $q/n$ at the beginning of the phase;
  \item each processor $v$ records an $n$-bit word $W(v)$, where $W_i(v)=1$ if and only if
        a pebble (possibly its own) passed through $v$ in round $i$ of the phase;
  \item if all words $W(v)$ are distinct, the processor with the lexicographically smallest word
        becomes the leader; otherwise a new phase starts.
\end{enumerate}

\noindent The algorithm relies on two invariants:
\begin{enumerate}
  \item every pebble disappears at the end of its phase, so the words depend only on the
        current phase;
  \item the words $W(v)$ are cyclic shifts of one word, since a pebble that passes $v_j$ in round
        $i$ passes $v_{j+k}$ in round $(i+k) \bmod n$. Hence every processor can compute all the
        words by shifting its own.
\end{enumerate}

\section{Correctness}

\begin{theorem}
If the algorithm terminates, exactly one processor is the leader.
\end{theorem}
\begin{proof}
At the end of a phase all processors see the same set of words, so either all of them start a new
phase or all of them stop. They stop only when the words are distinct, and then the smallest word
belongs to exactly one processor.
\end{proof}

\begin{theorem}
The algorithm terminates with probability $1$.
\end{theorem}
\begin{proof}
If exactly one pebble is sent, all words contain a single 1 at different positions, so the phase
succeeds. This happens with probability $q(1-q/n)^{n-1}>0$. Phases are independent, hence the
probability of $k$ consecutive failures tends to $0$.
\end{proof}

\section{Complexity}

\begin{theorem}[Itai, Rodeh, Theorem 2.2]
As $n\to\infty$, the expected number of phases tends to $(1-e^{-q})^{-1}$ and the expected number
of bits sent per processor tends to $q(1-e^{-q})^{-1}$.
\end{theorem}

The paper states this theorem without proof; below is our proof.

Each processor sends a pebble independently with probability $q/n$, and the pebble of $v_j$ passes
$v_0$ in round $(n-j) \bmod n$. Hence $W=W(v_0)$ consists of $n$ independent bits, each equal to $1$
with probability $q/n$. A phase fails if and only if $W$ equals its cyclic shift by some
$1\le k<n$; let $P_n$ be the probability of this event.

\begin{lemma}
If $W$ equals its shift by some $1\le k<n$, then for some divisor $r>1$ of $n$ the word $W$
consists of $r$ identical blocks.
\end{lemma}
\begin{proof}
Let $k$ be the smallest positive shift that does not change $W$. If $k$ did not divide $n$, then
$n=ak+b$ with $0<b<k$; shifts by $n$ and by $ak$ do not change $W$, so neither does the shift by
$b$, contradicting the choice of $k$. Thus $W$ consists of $r=n/k>1$ copies of a block of
length $k$.
\end{proof}

\begin{theorem}
For $n>q$
\[
\Bigl(1-\frac qn\Bigr)^{n}\ \le\ P_n\ \le\ \Bigl(1-\frac qn\Bigr)^{n}+\frac{2q^2}{\sqrt n},
\qquad\text{hence}\qquad P_n\to e^{-q}.
\]
\end{theorem}
\begin{proof}
The all-zero word, which occurs with probability $(1-q/n)^n$, equals all its shifts; this gives the
lower bound. If a failing word contains a $1$, then by the lemma it consists of $r$ identical blocks,
so this $1$ appears at the same position in all $r$ blocks. For a fixed $r$ there are $n/r$
positions, each giving $r$ ones with probability $(q/n)^r$, so by the union bound this case has
probability at most $\frac nr(q/n)^r\le q^2/n$. Since $n$ has at most $2\sqrt n$ divisors, summing
over $r$ gives at most $2q^2/\sqrt n$. Both bounds tend to $e^{-q}$.
\end{proof}

\begin{proof}[Proof of Theorem 2.2]
Phases are independent and each succeeds with probability $1-P_n$, so the number of phases is
geometric with expectation $1/(1-P_n)\to 1/(1-e^{-q})$. In every phase $n$ processors send a
pebble with probability $q/n$ each, so on average $q$ pebbles are sent. Each pebble makes $n$
steps, i.e.\ costs one bit per processor. Hence the expected number of bits per processor is
\[
q\cdot\mathbb E[\text{number of phases}]\ =\ \frac q{1-P_n}\ \longrightarrow\ \frac q{1-e^{-q}}.
\qedhere
\]
\end{proof}

\begin{thebibliography}{9}
\bibitem{ir} A.~Itai, M.~Rodeh, \emph{Symmetry breaking in distributed networks}, Information and
  Computation 88(1), 60--87 (1990).
\end{thebibliography}

\end{document}
